\item \textbf{Problema de repaso.} Dentro de la pared de una casa, una sección de tubería de agua caliente en forma de L consiste en tres partes: una pieza recta horizontal de $h = 28.0\text{ cm}$ de longitud; un codo; y una pieza vertical recta de $\ell = 134\text{ cm}$ de longitud (figura P6.6). Una trabe y un castillo mantienen fijos los extremos de esta sección de tubería de cobre. Encuentre la magnitud y dirección del desplazamiento del codo cuando hay flujo de agua, lo que eleva la temperatura de la tubería de 18.0 a 46.5 °C.
